\section{Discussion}
\label{sec:discussion}

During the development of BiteBound, several technical challenges were encountered, each requiring specific solutions to meet the set requirements.

As already mentioned in~\Cref{subsec:display}, the SPI transmission overhead posed a significant challenge to maintaining a stable 50 Hz frame rate. By implementing a hybrid rendering system using the dirty rectangles technique, a substantial reduction in payload size was achieved, allowing for faster transmission times.

Recursive algorithmic approaches could lead to stack overflow errors without setting maze size limits. Therefore, the iterative DFS implementation for maze generation was a necessary design choice to generate arbitrary maze sizes without limiting the system to an upper bound.

Sensor value noise and drift were expected at the beginning of the project. Implementing a deadzone threshold and an EMA filter provided a straightforward solution to ensure stable physical interactions while still allowing for raw telemetry data to be available for diagnostic purposes.

Testing the physics engine after implementation revealed that at high tilt angles, the ball could tunnel through maze walls due to its accelerated velocity. The implementation of sub-stepping during collision detection effectively mitigated this issue, ensuring that boundaries were respected even at high velocities, as described in~\Cref{subsec:physics_simulation_game_logic}.

During the architecture phase, the assumption was made that handling incoming dashboard commands and running the game loop on one core could cause significant frame drops. Early in development, the two-core possibility was explored, included in the initial design, and later implemented (see ~\Cref{sec:technical_concept}). This decision proved crucial for the modular system design and, consequently, the testability, maintainability, and scalability of the system; the dual-core synchronization strategy worked as expected.

As only one LCD display was available, the other team member had to acquire a separate ESP32-S3 board to test and run the hardware code. Though, the other board only had one core, which limited the ability to test the dual-core synchronization strategy.

The HiveMQ broker was found to be inconsistent when handling multiple subscriptions simultaneously, leading to lags and disconnects. Still, no other broker was chosen for further development, as this issue only appeared sporadically.